\chapter{Introduction}
\label{ch:introduction}
AutoPas is a particle simulation library, supporting molecular dynamics applications with a choice of particle containers, traversal algorithms and data layouts.
A time-stepping procedure advances the particle system until a predefined end time is reached.
Each time step requires computing force interactions between particles and updating the attributes of each individual particle.
\newline
\newline
A naive implementation that evaluates forces between all particle pairs quickly exceeds the realm of what is feasible for any sufficiently large simulation.
Consequently, AutoPas implements neighbor-search algorithms to restrict force evaluation to relevant nearby particles.
One such approach is the use of Verlet lists which store a list of neighbors for each particle, reducing the arithmetic work in each time step.
\newline
\newline
However, high-performance computing applications are also limited by memory access behavior.
In the case of Verlet lists, force interactions are calculated between a particle and each member in its respective neighbor list.
Critically, these particles, although close in the simulation domain, need not be close to each other in memory.
As a result, calculating force interactions with Verlet list may suffer from poor cache locality.
\newline
\newline
This bachelor's thesis investigates whether global spatial particle sorting can improve memory locality in Verlet lists.
The central question is whether improved locality can translate into runtime gains that compensate for the additional overhead introduced by sorting.
Furthermore, the thesis analyzes how different sorting resolutions and spatial ordering strategies behave across differently sized simulation scenarios.






